% Figure from MATH 556 notes, section 8. Compile with the notes preamble.
\begin{tikzpicture}[scale=2.0]
  % small ball around 0
  \fill[figB!8] (0,0) circle (0.3);
  \draw[-stealth, thin] (-1.35,0) -- (1.45,0) node[right, font=\scriptsize] {$x_1$};
  \draw[-stealth, thin] (0,-0.85) -- (0,1.7) node[above, font=\scriptsize] {$x_2$};
  % two lines through 0
  \draw[black!40, thin] (-0.85,-0.85) -- (1.35,1.35);
  \draw[black!40, thin] (-1.35,0.81) -- (1.35,-0.81);
  % the part of one line that stays in K
  \draw[figB, very thick] (-0.6,0.36) -- (0.6,-0.36);
  % the removed parabola
  \draw[figR, thick, domain=-1.25:1.25, samples=80] plot (\x,{\x*\x});
  \draw[figB, dashed, thin] (0,0) circle (0.3);
  \node[figB, font=\scriptsize, anchor=north east] at (-0.32,-0.06) {$B_r(0)$};
  % each line meets P once
  \fill[figR] (1,1) circle (0.85pt);
  \draw[figR, fill=white] (-0.6,0.36) circle (0.85pt);
  % the points (1/n, 1/n^2) inside the ball
  \foreach \n in {3,...,9} \fill[figR] ({1/\n},{1/(\n*\n)}) circle (0.55pt);
  \fill[black] (0,0) circle (0.8pt);
  \node[font=\scriptsize, anchor=north west, inner sep=1pt] at (0.03,-0.1) {$0$};
  \node[figR, font=\scriptsize, anchor=west] at (1.27,1.58) {$P = \{(s,s^2) : s \neq 0\}$};
  \node[font=\scriptsize, anchor=west] at (0.15,-0.75) {$K = \mathbb{R}^2 \setminus P$};
\end{tikzpicture}
